\section{\makebox[\widthof{\faBriefcase}][c]{\color{CVBlue}\faBriefcase}\ 实习经历}

\textbf{蚂蚁集团 · 客户体验与权益保障 · NLP实习算法工程师} \hfill 杭州 \quad 2026.06 -- 2026.09 \titlebreak
项目简述：智能客服拟人化改写系统（Qwen3-14B / 235B-A22B · verl · 40× H20）——针对“机器味重、易被识别后要求转人工”的痛点，把“更直接、更简洁、更类人”拆成可量化目标；参与数据构造、SFT / RL 训练、离线评测与线上 A/B 端到端 pipeline，与师兄完成对比实验设计并落地线上 A/B；已切 50\% 流量，质疑人工率 ↓16.7\% rel、平均回复长度 ↓23.4\%。
\begin{itemize}[nosep]
    \item \textbf{数据集构建（参与）：}针对“无现成智能 / 人工回复 pair”的难题，参与“逆向蒸馏 + 正向模拟”两段式构造，用 TF-IDF + 余弦相似度（阈值 0.8）去相近样本，以 LLM-judge 多维评测（底线合规 / 解决力对齐 / 按键满意度）清洗，产出 5,983 条多轮 + 6,000 条单轮样本。
    \item \textbf{评测体系与线上 A/B：}离线定义“混淆度 + 信息一致率”，线上接入“质疑人工率 / 转人工率 / 回复长度”，形成“数据 → 训练 → 评测 → 线上 AB → 回流”闭环；线上灰度按用户哈希分流并检查 SRM，质疑人工率 8.8\% → 7.3\%（↓16.7\% rel）、平均回复长度 ↓23.4\%。
    \item \textbf{RL 训练与消融：}将自动化评测能力构建为业务 reward 嵌入 RL 流程；针对 GRPO 仅序列级打分导致 token 分布振荡的问题，引入 token 级奖励（关键词 / 链接覆盖 + 长度 / 重复 / 分段合理性）+ seq-tanh 幅度调制；基于 verl 实现 GRPO baseline 并主导对比消融，拟人度持平（0.633 vs 0.638）下信息-效率 95.8\% → 98.4\%（+2.6 pt），优于纯 SFT（96.8\%）。
\end{itemize}
